% Part:sets-functions-relations
% Chapter: size-of-sets
% Section: pairing

\documentclass[../../../include/open-logic-section]{subfiles}

\begin{document}

\olfileid{sfr}{siz}{pai}
\olsection{જોડ-નિરૂપણ વિધેયો અને સંકેતસંખ્યાઓ}

\begin{explain}
કૅન્ટૉરની આડીઅવળી રીત $\Nat^n$ની ગણનીયતા દૃષ્ટિગોચર
બનાવે છે. પરંતુ $\Nat^2$ દર્શાવતી આપણી સરણિ પર
ધ્યાન આપીએ. સરણિમાં આડીઅવળી રેખાને અનુસરીને સ્થાનો
ગણતાં તપાસી શકીએ કે $\tuple{1,2}$ સાથે સંખ્યા~$7$
જોડાય છે. જોકે, તેની વધુ સીધી ગણતરી કરી શકીએ તો
સારું થાય. એટલે કે આડીઅવળી પરિગણનાનું
\emph{પ્રતિવિધેય} $g\colon \Nat^2 \to \Nat$
આપણી પાસે હોય તો સારું, જેથી
\[
g(\tuple{0,0}) = 0, \;
g(\tuple{0,1}) = 1, \;
g(\tuple{1,0}) = 2, \; \dots,
g(\tuple{1,2}) = 7, \; \dots
\]
થાય. તેનાથી પરિગણનામાં $\tuple{n, m}$ ચોક્કસ
ક્યાં આવશે તેની ગણતરી કરી શકીએ.

હકીકતમાં, બે નોંધોના આધારે $g$ને સીધું વ્યાખ્યાયિત
કરી શકીએ. પ્રથમ: $n$મી હાર અને $m$મા સ્તંભમાં
કિંમત~$v$ હોય, તો $(n+1)$મી હાર અને $(m-1)$મા
સ્તંભમાં કિંમત $v + 1$ હોય છે. બીજી: પરિગણનાની
પ્રથમ હારમાં ત્રિકોણીય સંખ્યાઓ આવે છે, જે
$0$, $1$, $3$, $6$ વગેરેથી શરૂ થાય છે.
$k$મી ત્રિકોણીય સંખ્યા એ $< k$ હોય તેવી પ્રાકૃતિક
સંખ્યાઓનો સરવાળો છે, જેની ગણતરી $k(k+1)/2$થી
કરી શકાય. આ બંને નોંધોને ભેગી કરીને નીચેનું વિધેય લો:
\[
  g(n,m) = \frac{(n+m+1)(n+m)}{2} + n
\]
આપણે ઘણી વાર $g(\tuple{n, m})$ને બદલે માત્ર
$g(n, m)$ લખીએ છીએ, કારણ કે તે વાંચવામાં સરળ છે.
આ વિધેય કહે છે કે પહેલાં $(n+m)^\text{મી}$
ત્રિકોણીય સંખ્યા નક્કી કરો અને પછી તેમાં $n$ ઉમેરો.
તે સરણિને આપણી ઇચ્છા મુજબ જ ભરે છે.
ખાસ કરીને, જોડ $\tuple{1, 2}$ને
$\frac{4 \times 3}{2} +
1 = 7$ પર મોકલવામાં આવે છે.

આ વિધેય $g$ એ જોડોના ગણની પરિગણનાનું
\emph{પ્રતિવિધેય} છે. આવાં વિધેયોને
\emph{જોડ-નિરૂપણ વિધેયો} કહે છે.
\end{explain}

\begin{defn}[જોડ-નિરૂપણ વિધેય]
વિધેય $f\colon A \times B \to \Nat$ અંકગણિતીય
\emph{જોડ-નિરૂપણ વિધેય} છે જો $f$ એક-એક હોય.
આપણે એમ પણ કહીએ છીએ કે $f$ એ $A \times B$ને
\emph{સંકેતબદ્ધ કરે છે}, અને $f(x,y)$ એ
$\tuple{x,y}$ની \emph{સંકેતસંખ્યા} છે.
\end{defn}

\begin{explain}
જોડ-નિરૂપણ વિધેયો વડે, દાખલા તરીકે, પ્રાકૃતિક
સંખ્યાઓની જોડોને સંકેતબદ્ધ કરી શકીએ; બીજા શબ્દોમાં,
ઘટકોની દરેક \emph{જોડ}ને \emph{એક જ} સંખ્યા વડે
દર્શાવી શકીએ. જોડ-નિરૂપણ વિધેયના પ્રતિવિધેય વડે
સંખ્યાનો \emph{સંકેત ઉકેલી શકીએ}, એટલે કે તે
કઈ જોડ દર્શાવે છે તે શોધી શકીએ.
\end{explain}

\begin{prob}
બધી અનૃણ સંમેય સંખ્યાઓના ગણની પરિગણના આપો.
\end{prob}

\begin{prob}
બતાવો કે $\Rat$ એ !!{enumerable} છે.
યાદ કરો કે કોઈ પણ સંમેય સંખ્યાને $z/m$ અપૂર્ણાંક
રૂપે લખી શકાય, જ્યાં $z \in \Int$ અને $m \in \Nat^+$.
\end{prob}

\begin{prob}
$\Bin^*$ની પરિગણના વ્યાખ્યાયિત કરો.
\end{prob}

\begin{prob}
તમારા પ્રારંભિક તર્કશાસ્ત્રના અભ્યાસમાંથી યાદ કરો
કે દરેક સંભવિત સત્ય કોષ્ટક એક સત્ય વિધેય વ્યક્ત
કરે છે. બીજા શબ્દોમાં, સત્ય વિધેયો એ કોઈક~$k$
માટે $\Bin^k \to \Bin$નાં બધાં વિધેયો છે.
સાબિત કરો કે બધાં સત્ય વિધેયોનો ગણ ગણનીય છે.
\end{prob}

\begin{prob}
બતાવો કે કોઈ પણ અનંત !!{enumerable} ગણના
બધા સાન્ત ઉપગણોનો ગણ !!{enumerable} છે.
\end{prob}

\begin{prob}
$\Nat$નો ઉપગણ \emph{સહસાન્ત} કહેવાય જો અને માત્ર
જો તે $\Nat$ના કોઈ સાન્ત ઉપગણનો પૂરક હોય; એટલે કે
$A \subseteq \Nat$ સહસાન્ત હોય જો અને માત્ર જો
$\Nat\setminus A$ સાન્ત હોય. ધારો કે $I$ એવો ગણ
છે જેના !!{element}s ચોક્કસપણે $\Nat$ના સાન્ત
અને સહસાન્ત ઉપગણો છે. બતાવો કે $I$ એ !!{enumerable} છે.
\end{prob}
\sourcecorrection{OLSIZ-002}{મૂળની \texttt{pairing.tex}, પંક્તિઓ 91--97માં
``a finite set \(\Nat\)'' એવો વ્યાકરણદોષ છે. પછીનું સમીકરણ અભિપ્રાય
નક્કી કરે છે: સહસાન્ત ઉપગણ એ \(\Nat\)ના કોઈ સાન્ત ઉપગણનો પૂરક છે.}

\begin{prob}
બતાવો કે !!{enumerable} ગણોનો !!{enumerable}
યોગગણ !!{enumerable} છે. એટલે કે જ્યારે
$A_1$, $A_2$, \dots{} ગણો હોય અને દરેક $A_i$
એ !!{enumerable} હોય, ત્યારે તે બધાનો યોગગણ
$\bigcup_{i=1}^\infty
A_i$ પણ !!{enumerable} છે. [નોંધ: આ અઘરું છે!]
\end{prob}

\begin{prob}
ધારો કે $f \colon A \times B \to \Nat$ કોઈ પણ
જોડ-નિરૂપણ વિધેય છે. બતાવો કે $f$નું પ્રતિવિધેય
$A \times B$ની પરિગણના છે.
\end{prob}

\begin{prob}
$\Nat^3$ને સંકેતબદ્ધ કરતું વિધેય નિર્દિષ્ટ કરો.
\end{prob}

\end{document}
